\section{Small-scale geometry and the outermost scale}\label{sec:geometry}

Throughout this section, $(M^4,g(t))$, $0\leq t<T$, is the closed
Ricci flow under consideration and $|R|\leq C_R$. Constants may depend
on this flow, $T$, and a specified finite differentiation order. They are
independent of the small scales subsequently chosen. For a tensor $A$,
write $|A|_{k,r}=\sum_{j=0}^k r^j|\nabla^j A|$; on coordinate annuli
the analogous notation uses Cartesian derivatives. We first establish
the geometric input needed in the later analytic argument. The
bounded-scalar-curvature estimates in
\cite{BamlerZhangI,BamlerZhangII} give background for these estimates;
the scale selection, tree organization, and motion assertions needed
here are included below.

\begin{definition}[Spatial curvature radius]\label{geom:radius}
Fix a small cap $l_0>0$. At a material point $x\in M$ set
\[
 l(x,t)=\sup\{0<r\leq l_0:
       \sup_{B_t(x,r)}|\Rm|\leq r^{-2}\}.
\]
The definition uses the whole spatial ball, not just curvature at its
center. The function $l$ is positive and continuous before $T$, and
$|l(x,t)-l(y,t)|\leq d_t(x,y)$. Indeed a ball of radius $r-d_t(x,y)$
at $y$ is contained in the radius-$r$ ball at $x$, and its allowed
curvature bound is larger. Continuity in time follows by compactness
and smooth dependence of the metric, using radii slightly smaller or
larger than the supremum.
\end{definition}

\begin{theorem}[Scale estimates and degeneration]\label{thm:degeneration}
For a sufficiently small cap and all sufficiently late times, the
following conclusions hold.

\emph{Scale and energy bounds.} One has
\begin{align}
 cr^4\leq\Vol_t B_t(x,r)&\leq Cr^4,
       &\int_M|\Rm|^2\,d\mu_t&\leq C,\label{geom:volume-energy}\\
 |\nabla^k\Rm|(x,t)&\leq C_k l(x,t)^{-2-k},
       &|\nabla^k\Ric|(x,t)&\leq C_k l(x,t)^{-1-k}.
       \label{eq:ricci-improvement}
\end{align}
Here $0<r\leq l_0$ and each fixed $k$ is allowed. The set
$\{l<r\}$ is covered by a uniformly bounded number of balls of radius
$Cr$. Distances converge uniformly to a compact metric quotient $Z$
of the material manifold. Away from finitely many concentration points
the metric convergence is smooth.

\emph{Small-scale limits and genuine singularities.}
At arbitrary sequences of centers and length scales $h_i\to0$, with
times tending to $T$, a subsequence of the rescaled slices converges
to a complete noncollapsed Ricci-flat orbifold, smoothly away from
finitely many genuine orbifold points. Each such point has a nontrivial
finite group acting freely on $S^3$. A nonflat limit has a single
Euclidean quotient end with nontrivial group. A flat limit is a global
Euclidean quotient and has at most one exceptional point.

\emph{Finite trees.} Within each concentrating cluster, maximal
inequivalent nonflat limits form a finite
rooted tree with a unique outermost entry. The joining regions have the
quantitative charts of Lemma~\ref{geom:neck} below.
\end{theorem}

Sections~\ref{geom:local-estimates}--\ref{geom:energy-comparison}
prove the scale and energy estimates and convergence of distances.
Sections~\ref{geom:slice-limits}--\ref{geom:puncture-removal} establish
the small-scale limits and identify their exceptional points as genuine
quotient singularities. Section~\ref{geom:trees-root} completes terminal
regular convergence and constructs the finite trees. The remaining
subsections use these conclusions to compare the outermost scale in time.

\subsection{Heat localization and two-sided persistence of the spatial radius}
\label{geom:local-estimates}

The spatial curvature radius does not by itself control a time cylinder.
Our first goal is to obtain that control in both time directions, and
then to improve the Ricci bound by one power of the spatial radius.
We give the heat estimates first so that the curvature-radius argument
does not assume curvature-dependent cutoff bounds. Set $n=4$ for the
applications, retaining $n$ temporarily in the heat formulas. If
$u=(4\pi\tau)^{-n/2}e^{-f}$ is a positive unit-mass conjugate heat
density, $\partial_\tau=-\partial_t$ and
$\partial_\tau u=\Delta u-Ru$, put
\[
 v=\{\tau(2\Delta f-|\nabla f|^2+R)+f-n\}u.
\]
A direct computation gives
\begin{equation}\label{geom:entropy-identity}
 (\partial_\tau-\Delta+R)v
 =-2\tau|\Ric+\nabla^2f-g/(2\tau)|^2u.
\end{equation}
For completeness, after factoring out $u$ the operator is
$D=\partial_\tau-\Delta+2\nabla f\cdot\nabla$. The identities used are
\begin{align*}
 f_\tau&=\Delta f-|\nabla f|^2+R-n/(2\tau),\\
 D\Delta f&=-2|\nabla^2 f|^2-2\Ric(\nabla f,\nabla f)
                  +\Delta R-2\Ric:\nabla^2f,\\
 D|\nabla f|^2&=2\nabla f\cdot\nabla R
                   -2|\nabla^2 f|^2-4\Ric(\nabla f,\nabla f),\\
 DR&=-2\Delta R-2|\Ric|^2+2\nabla f\cdot\nabla R.
\end{align*}
The metric term in $\partial_\tau\Delta f$ is
$-2\Ric:\nabla^2f$; the contracted Bianchi identity cancels the
first-derivative terms. Substitution gives
Equation~\eqref{geom:entropy-identity}, including its time direction.
This is Perelman's entropy calculation
\cite[Proposition~9.1]{PerelmanEntropy}.

Consequently $\int v\,d\mu_t$ is nondecreasing in forward time.
At time zero it has a lower bound uniform for $0<\tau\leq T+1$:
writing $\phi=\sqrt u$ and $\int\phi^2=1$, its expression is
\[
 \tau\int(4|\nabla\phi|^2+R\phi^2)
 -\int\phi^2\log\phi^2-\frac n2\log(4\pi\tau)-n.
\]
Jensen's inequality followed by the Sobolev inequality of the fixed
initial metric gives this bound, including as $\tau\downarrow0$ by
optimizing the coefficient of $\|\nabla\phi\|_2^2$. Solve the
conjugate equation backwards from an arbitrary smooth positive
terminal density. Monotonicity then gives the same logarithmic Sobolev
inequality at every time for $0<\tau\leq1$. Positive approximation
allows Lipschitz or nonnegative test functions.

Here are two consequences with their normalizations. If a small ball
has arbitrarily small ratio $\Vol(B_r)/r^n$, choose a dyadic descendant
whose inner half has a definite fraction of its volume and whose ratio
is still small. Such a descendant exists: otherwise successive ratios
decrease geometrically to zero, contradicting the Euclidean ratio at a
smooth center. A tent function supported in this descendant, equal to
one on its inner half, has normalized Dirichlet energy $O(r^{-2})$.
Jensen gives $\int\phi^2\log\phi^2\geq-\log\Vol(B_r)$.
The logarithmic Sobolev inequality with $\tau=r^2$ therefore gives a
uniform positive lower volume ratio. Also
$\int\phi^2\log\phi^2\geq-2\log\|\phi\|_1$ for
$\|\phi\|_2=1$. Optimization in $\tau$, using $\tau=1$ when
necessary, gives
\begin{equation}\label{geom:nash}
 \|\phi\|_2^{2+4/n}
 \leq C(\|\nabla\phi\|_2^2+\|\phi\|_2^2)\|\phi\|_1^{4/n}.
\end{equation}
For both forward ordinary heat and backwards conjugate heat, the
$L^1$ norm is bounded by its initial norm times $e^{Ch}$, and the
derivative of the squared $L^2$ norm in the diffusion direction is at
most $-2\|\nabla u\|_2^2+C\|u\|_2^2$. This uses only
$\partial_t d\mu_t=-R\,d\mu_t$. Equation~\eqref{geom:nash} and
integration of the resulting scalar differential inequality give
$\|P_h\|_{1\to2}\leq Ch^{-n/4}$ for $h\leq1$. Duality with the
conjugate equation and composition at the midpoint give the heat
kernel bound $K\leq Ch^{-n/2}$.

For a conjugate kernel with terminal pole $p$,
Equation~\eqref{geom:entropy-identity} also gives $v\leq0$; compare
\cite[Corollaries~9.3--9.4]{PerelmanEntropy}.
One may verify the terminal sign without a uniform curvature bound
near $T$. Fix the terminal time first, when the flow is smooth, and a
positive forward heat test $\zeta$. Integration by parts gives
\[
 \int v\zeta
 =\partial_\tau\left(\tau\int(f-n/2)\zeta u\right)
       -2\tau\int u\Delta\zeta.
\]
The kernel upper bound bounds the integral in parentheses from below
after division by $\tau$. For its upper bound, apply Jensen relative
to $(4\pi\tau)^{-n/2}e^{-d_*^2/(4\tau)}$, where $d_*^2$ is a smooth
function agreeing with squared terminal distance near $p$ and has
$p$ as its unique zero. Its mass tends to $1$, and the heat equation,
tested against $d_*^2\zeta$, gives
$\tau^{-1}\int d_*^2\zeta u\to2n\zeta(p)$. Hence
$\limsup\int(f-n/2)\zeta u\leq0$. Along a sequence approaching zero
the displayed derivative has nonpositive liminf. At this fixed smooth
terminal time, $2\tau\int u\Delta\zeta\to0$, so the same is true of
$\int v\zeta$. Equation~\eqref{geom:entropy-identity} makes
$\int v\zeta$ nonincreasing in $\tau$; hence it is nonpositive at
every earlier time. Arbitrary positive forward heat tests now give
$v\leq0$. Equivalently,
\[
 2\tau f_\tau+\tau|\nabla f|^2-\tau R+f\leq0.
\]
Integration along the constant material path at $p$ gives
$f(p,\tau)\leq C\tau$. The initial term in this integration can
also be obtained from the preceding second moment: some point within
$O(\sqrt\tau)$ of $p$ has bounded $f$; moving from that point to $p$
over a time interval of size $\tau$ costs $O(\sqrt\tau)$ in the
inequality for $\sqrt\tau f$, by completing the gradient square.
Letting the starting time tend to zero removes this cost. Thus
\begin{equation}\label{geom:kernel-diagonal}
 K(p,t;p,s)\geq c(t-s)^{-n/2}.
\end{equation}

For a forward heat solution $|U|\leq1$ over elapsed time $h\leq1$,
the cancellation in the Bochner formula gives
\[
 (\partial_t-\Delta)|\nabla U|^2=-2|\nabla^2U|^2,
 \qquad |\nabla U|\leq Ch^{-1/2}.
\]
The gradient bound follows by applying the maximum principle to
$h|\nabla U|^2+CU^2$. To obtain a time derivative bound observe
$ (\partial_t-\Delta)\Delta U=2\Ric:\nabla^2U$ and
$ (\partial_t-\Delta)R=2|\Ric|^2$. For either sign, apply the
maximum principle to
$Q=h^2(\pm\Delta U+A|\nabla U|^2-AR)$, with fixed large $A$.
The negative terms $-2Ah^2(|\nabla^2U|^2+|\Ric|^2)$ absorb
$2h^2\Ric:\nabla^2U$. The remaining negative Hessian square
absorbs $2h|\Delta U|$; the already proved gradient bound controls
the other elapsed-time terms. At a positive maximum with $Q\geq C h$
the resulting upper bound is $C-cQ^2/h^2$, which rules out a first
crossing of a sufficiently large multiple of $h$. Therefore
\begin{equation}\label{geom:heat-derivatives}
 |\nabla U|\leq Ch^{-1/2},\qquad |\partial_tU|=|\Delta U|\leq C/h.
\end{equation}

Take a forward kernel from $(x,t-r^2)$ and multiply it by $r^n$.
On a short slab about $t$ it is bounded, has value at least $c$ at
$(x,t)$, has spatial gradient $C/r$, time derivative $C/r^2$, and
integral at most $Cr^n$. It is therefore bounded below on a ball of
radius $cr$ at time $t$, proving the upper volume bound. Each
connected component of a fixed positive superlevel has diameter
$O(r)$: a connecting path of greater diameter contains arbitrarily
many disjoint balls of radius $c'r$, where the function is bounded
below; lower volume bounds and the integral estimate bound their
number. The same function, using its time derivative bound and
slightly nested levels, proves
\begin{equation}\label{geom:rough-ball-transfer}
 B_{t_1}(x,cr)\subset B_{t_2}(x,Cr)
       \quad\text{if }|t_1-t_2|\leq c'r^2,
\end{equation}
in either time direction. Indeed every initial radial path in the
smaller ball remains in a positive superlevel at the second time, so
its endpoint remains in the same component as $x$.

In particular, for small $l$ there is a nonnegative Lipschitz cutoff
on $[t-cl^2,t]$, bounded by a fixed constant, with
\begin{equation}\label{geom:heat-cutoff}
 |\nabla\chi|\leq C/l,\quad
 c/l^2\leq\partial_s\chi\leq C/l^2\ \text{ on }\{\chi>0\},
 \quad\chi(x,s)\geq c,\quad\supp\chi\subset B_t(x,l).
\end{equation}
Choose the kernel radius to be a small fixed fraction of $l$, subtract
a suitable level, add $C'(s-t)/r^2$, and take the positive part. Keep
the final superlevel component containing $x$. Choosing $C'$ larger
than the kernel time-derivative constant makes the positive domains
monotone in $s$, so extension by zero outside this component is
legitimate. The component diameter bound puts its support in
$B_t(x,l)$. All constants are independent of the cap $l_0$.

We now bootstrap, rather than assume, the time control needed to use
a spatial curvature radius. Fix a tensor-index convention and
bootstrap
\begin{equation}\label{geom:bootstrap}
 |\partial_t\Rm|+C_n|\Ric||\Rm|\leq A l^{-3}.
\end{equation}
The left side controls the time derivative of the curvature norm.
Assume Equation~\eqref{geom:bootstrap} at preceding times, set
$l=l(x,t)$, and use Equation~\eqref{geom:heat-cutoff}. At time $t$,
$\max\chi^2|\Rm|\leq C l^{-2}$. Choose $B>C$. If this maximum
had been larger than $Bl^{-2}$, take its last crossing down to that
value, and let $y$ be a maximizing point at the crossing. Put
$\chi_y=\chi(y,s)>0$. The maximum and the cutoff gradient bound give,
for $d_s(y,z)\leq c\chi_y l$,
\[
 \chi(z,s)\geq\chi_y/2,\qquad
 |\Rm|(z,s)\leq4B/(\chi_y^2 l^2).
\]
Thus, taking $c_B\leq\min(c,(4B)^{-1/2})$,
\begin{equation}\label{geom:no-circularity}
 l(y,s)\geq c_B\chi_y l.
\end{equation}
This inference uses no curvature derivative estimate. At the crossing
the positive cutoff contribution to $\partial_s(\chi^2|\Rm|)$ is
at least $cB/(\chi_y l^4)$, whereas its possible negative curvature
contribution is at most $C_BA/(\chi_y l^3)$ by the bootstrap.
Choosing $A l_0$ sufficiently small makes the derivative strictly
positive. The upper one-sided derivative of the spatial maximum
cannot then permit the crossing. Approximation of the positive-part
cutoff, or the corresponding one-sided maximum argument on its
positive set, gives the same conclusion. Consequently curvature is
bounded by $Cl^{-2}$ on a definite smaller backward parabolic ball
at $(x,t)$. Its moving spatial balls are contained in the positive
cutoff region by $\chi(x,s)\geq c$ and the gradient bound.

On this parabolic ball the usual differentiated curvature argument
can now be applied without circularity. The connection variation is
\[
 \partial_t\Gamma^k_{ij}
 =-g^{km}(\nabla_i\Ric_{jm}+\nabla_j\Ric_{im}-\nabla_m\Ric_{ij}).
\]
Skew differentiation, the second Bianchi identity, and commuting
derivatives give the heat equation for $\Rm$ and its differentiated
equations, with products
$\nabla^i\Rm*\nabla^{k-i}\Rm$. At unit scale, after bounds through
order $k-1$, the quantity
$(D+|\nabla^{k-1}\Rm|^2)|\nabla^k\Rm|^2$, with large fixed $D$,
satisfies a differential inequality bounded above by a constant minus
a positive multiple of its square. Cauchy's inequality absorbs the
cross-gradient terms. Cutoffs on smaller parabolic balls, whose
distance and Laplacian bounds now follow from the curvature bound,
prove the derivative estimate inductively; distance cutoffs are
interpreted by shortened radial barriers at cut points. This is the
local derivative argument of \cite{Shi1989}.

The coordinate estimates used here have a uniform radius. Normalize
a smaller controlled curvature ball by the length $l$. On a fixed
ball in these units, sectional curvature is bounded above and below,
and every sufficiently small contained ball has volume at least
$c r^4$. The local injectivity estimate of Cheeger--Gromov--Taylor
\cite[Theorem~4.3]{CheegerGromovTaylor1982} therefore gives a positive
uniform injectivity radius on a further smaller ball. To check its
local hypotheses, reserve a fixed outer ball within the curvature
bound and choose all comparison radii below its conjugate-radius
bound. The upper sectional bound supplies that bound; the lower
sectional bound controls the volume of the pulled-back tangent ball
in the theorem's loop estimate, while the lower ball-volume bound
controls its denominator. Closedness supplies compactness of the
smaller closed balls. No curvature bound outside this controlled
region is used. Returning to original units gives normal radius
$c l$. Differentiating radial Jacobi equations then gives the higher
coordinate bounds from the covariant curvature bounds.

At this stage the crude estimates
$|\nabla^k\Ric|\leq C_k l^{-2-k}$ follow by contraction from the
curvature estimates. These already suffice, after dilation to unit
scale, to integrate the metric and connection equations in a fixed
initial normal chart for a time of order $l^2$. Consequently the
chart coefficients and all needed derivatives stay uniformly
controlled, and smooth spacetime cutoffs with the usual scaled
bounds are available before the improved Ricci estimate is proved.

The lower-index Ricci tensor satisfies
\begin{equation}\label{geom:ricci-equation}
 \partial_t\Ric=\Delta\Ric+2\Rm(\Ric)-2\Ric^2,
       \qquad \Rm(g)=\Ric.
\end{equation}
Test the scalar equation against a scaled smooth spacetime cutoff on
the smaller parabolic ball. Such cutoffs are now available in the
controlled coordinates. Integration by parts, $|R|\leq C_R$, and
volume $O(l^4)$ over time $O(l^2)$ give
\begin{equation}\label{geom:scalar-energy}
 \int_{t-cl^2}^{t}\int_{B_s(x,cl)}|\Ric|^2\,d\mu_s\,ds
       \leq Cl^4.
\end{equation}
For example, derivatives of the cutoff cost $l^{-2}$, so its scalar
terms cost $l^{-2}l^4l^2= l^4$; time integration also includes the
bounded $R^2$ measure term. Rescaling lengths by $l^{-1}$ makes the
squared spacetime Ricci norm $O(l^2)$. Equation~\eqref{geom:ricci-equation}
is a homogeneous linear tensor heat equation once its bounded
coefficients are fixed. Interior energy estimates followed by local
Sobolev iteration and differentiated interior estimates give
$|\nabla^k\Ric|\leq C_k l^{-1-k}$ in original units.
This is the local scalar-to-Ricci improvement of
\cite[arXiv version, Lemma~6.1]{BamlerZhangI}; the preceding argument
has verified its bounded-curvature cylinder and cutoff inputs from
the spatial radius. The uniform coefficients used in these interior
estimates are precisely those obtained from the crude curvature bounds above;
the improved Ricci estimates are their conclusion.

Curvature variation now bounds the left side of
Equation~\eqref{geom:bootstrap} by $A_0 l^{-3}$, with $A_0$
independent of $A,l_0$. Choose $A>A_0$, and then $l_0$ small enough
for the crossing argument and for the bootstrap to hold on a fixed
initial compact time interval. A first-failure argument proves it for
all later times. This proves Equation~\eqref{eq:ricci-improvement}.
For forward persistence, use the same heat function with a negative
drift, so that $-C/l^2\leq\partial_s\chi\leq-c/l^2$ on its positive
set and its support stays in the initial radius-$l$ ball. At the
first upward crossing of $\max\chi^2|\Rm|=Bl^{-2}$, the cutoff
gradient again gives $l(y,s)\geq c_B\chi_y l$. The negative cutoff
contribution is at most $-cB/(\chi_y l^4)$ while the curvature
contribution is at most $C_BA/(\chi_y l^3)$. Their sum is negative
for the chosen cap, contradicting the crossing. The cutoff is
positive near the initial center throughout a fixed shorter time
slab, proving the forward bounded-curvature neighborhood. We have
proved both directions of spatial-scale persistence and the derivative
bounds in Equation~\eqref{eq:ricci-improvement}; the two-sided volume
bounds were established by heat localization. The remaining global
input for compactness is the slice energy bound.

\subsection{Energy, small bad sets, and volume comparison}
\label{geom:energy-comparison}

The local estimates now give definite curvature energy at each
concentration scale. To bound how many such regions can occur, we first
control the total slice energy and then derive a summable volume-comparison
error. The latter will distinguish genuine quotient singularities from
smooth points in a limit.
Choose $C_0$ so $1\leq R+C_0\leq C$. Completing the gradient square
in the quotient evolution, as in Simon's integral-curvature argument
\cite[Section~3, equation~(3.2) and Theorems~3.5--3.6]{SimonIntegral2020}, gives
\[
 (\partial_t-\Delta)\frac{|\Ric|^2}{R+C_0}
 \leq C|\Rm||\Ric|^2-c|\Ric|^4.
\]
If $F=\int|\Ric|^2/(R+C_0)$, the volume derivative and Young's
inequality imply
\[
 F'\leq C\int|\Rm|^2-c'\int|\Ric|^4+CF.
\]
In four dimensions the Chern--Gauss--Bonnet formula \cite{Chern1944}
in the full-tensor normalization of
\cite[Equation~(3.6)]{SimonIntegral2020} gives
$\int|\Rm|^2\leq C+4\int|\Ric|^2\leq C+CF$. No orientability
assumption is needed: for a nonorientable $M$ apply the identity to
its orientation double cover with the pulled-back flow and divide
the integral identity by two. Dropping the negative quartic
term and applying Gronwall bounds $F$ and then the slice curvature
energy, proving Equation~\eqref{geom:volume-energy}.

There is a fixed energy $e_0>0$ associated to a point-picked
curvature ball. Indeed maximize curvature times squared distance to
the boundary of a slightly larger ball. At the maximizing point,
curvature $Q$ is bounded by $CQ$ on a ball of radius $cQ^{-1/2}$.
The spatial derivative estimate makes curvature at least $Q/2$ on a
smaller ball; lower volume then gives $\int|\Rm|^2\geq e_0$.
If $l(x)<r$, curvature reaches at least $r^{-2}$ within distance $r$
of $x$ by the definition of $l$. The preceding point selection
therefore produces an energy ball inside $B(x,Cr)$. A maximal
disjoint collection of these larger radius-$Cr$ balls has at most
$C/e_0$ members by the slice energy bound; enlarging them covers
$\{l<r\}$. Thus
\begin{equation}\label{geom:bad-volume}
 \Vol_t\{l<r\}\leq Cr^4.
\end{equation}
On $B(x,h)\cap\{l\geq h\}$, the cubic Ricci integral is at most
$Ch^{-3}h^4$. On $2^{-j-1}h\leq l<2^{-j}h$ it is at most
$C(2^{-j}h)^{-3}(2^{-j}h)^4$. Summing gives
\begin{equation}\label{geom:ricci-L3}
 \int_{B_t(x,h)}|\Ric|^3\,d\mu_t\leq Ch.
\end{equation}

Distances have a uniform terminal limit. For $|t-s|\leq cr^2$,
cover the bad set at the first time by the bounded family of
radius-$Cr$ balls, omitting from a minimizing path its visits to
slightly enlarged balls. Replace first-to-last visits when a ball is
revisited. There are uniformly many replacement endpoints, with
total replacement cost $O(r)$ at the other time by
Equation~\eqref{geom:rough-ball-transfer}. On the remaining path the
curvature scale is $\geq cr$ throughout the interval by forward and
backward persistence. Its length changes by a factor at most $1+Cr$
since $|\Ric|\leq C/r$. Comparing also in reverse gives
\[
 |d_t(x,y)-d_s(x,y)|\leq Cr(1+d_s(x,y)+d_t(x,y)).
\]
First compare with a fixed late time to bound diameter, then take
$r$ a sufficiently large constant times $\sqrt{|t-s|}$. This proves
uniform convergence to a pseudodistance $d_T$ on $M$ and defines its
compact metric quotient $Z$. For scalar-bound distance distortion and
terminal metric convergence, see also
\cite[Theorem~1.1 and Corollary~1.2]{BamlerZhangII}.
Volume forms at two times differ by factors $e^{\pm C_R|t-s|}$.

We will need a comparison estimate stronger than mere volume bounds.
In polar coordinates about a smooth point, before the cut radius,
let $P$ be the polar Jacobian and
$m=(\partial_r\log P-3/r)_+$. The Riccati inequality gives, on
$\{m>0\}$,
\[
 m'+2m/r+m^2/3\leq|\Ric|.
\]
Since $P'=(m+3/r)P$ there,
\[
 (m^5P)'\leq5|\Ric|m^4P-7m^5P/r-(2/3)m^6P.
\]
Integrate to cuts and use Young's inequality on the first term.
The boundary contributions have the correct sign, with $m=0$ at
entrance into a positivity interval. Setting $P=0$ after cuts gives
\begin{equation}\label{geom:polar-defect}
 \int_{B(x,h)}m^6\leq C\int_{B(x,h)}|\Ric|^3\leq Ch.
\end{equation}
Furthermore $(P/r^3)'\leq mP/r^3$. Integrating between comparable
radii yields
\begin{equation}\label{geom:almost-bishop}
 \frac{\Vol B(x,h)}{h^4}
 \leq\frac{\Vol B(x,h/2)}{(h/2)^4}+Ch^{1/2}.
\end{equation}
To see the exponent, the error is bounded by
$Ch^{-3}\int m\leq Ch^{-3}(Ch)^{1/6}(Ch^4)^{5/6}=Ch^{1/2}$.
The same argument holds for any fixed comparable pair of radii.
Errors sum to $Ch^{1/2}$ over all dyadic descendants. Polar
integration including the nonpositive cut measure also gives weakly
\begin{equation}\label{geom:radial-laplacian}
 \Delta(d^2/2)\leq4+dm.
\end{equation}

The scale, energy, and bad-set bounds of Theorem~\ref{thm:degeneration}
are now established, as is uniform convergence of distances. We next
identify the geometry on every shrinking spatial scale; smooth terminal
convergence away from a fixed finite set will then follow by persistence.

\subsection{Slice limits, quotient cones, and quantitative necks}
\label{geom:slice-limits}

We next turn these estimates into a description of every small-scale
slice limit. The essential conclusions are smooth Ricci-flat convergence
off a finite set and quantitative quotient coordinates on the intervening
annuli. Puncture removal and the nontriviality of the quotient groups
will be established after those coordinates are available.
This is the part of the argument corresponding to classical Einstein
orbifold compactness and bubbling \cite{Anderson1989,Bando1990,
BandoCorrection1990}; the estimates below include the Ricci error
present in the flowing slices.
Consider $h_i^{-2}g(t_i)$, based at arbitrary $x_i$, where
$h_i\to0$ and $t_i\uparrow T$. Two-sided volume bounds uniformly
bound the number of disjoint radius-$\rho$ balls in any bounded ball,
giving pointed subsequential compactness as length spaces.
Completeness follows by exhausting with successively larger balls.
The cover of $\{l<\rho h_i\}$ has uniformly bounded cardinality.
Letting $\rho\downarrow0$ diagonally identifies a finite exceptional
set $S$, with $l/h_i$ bounded below on its regular compact complement.
The coordinate and derivative bounds give smooth convergence there.
In the rescaled metrics,
\[
 |\nabla^k\Ric|\leq C_kh_i(l/h_i)^{-1-k}\longrightarrow0
\]
on these compacts, so the limit is Ricci-flat.

Smooth convergence includes embeddings of relatively compact regular
subregions. Take finitely many normal charts with margins and extract
their transition maps together with their metrics. The local inverse
identifications differ smoothly by quantities tending to zero.
In convex limiting balls combine them by a smooth partition and the
unique local squared-distance barycenter. The resulting map is locally
invertible on smaller charts, covers the desired region by the chart
margins, and is injective: two points whose images converge to the same
point have distance tending to zero, hence lie in one chart where
injectivity holds. Construct on a larger compact region before
restricting. This supplies actual embeddings for the later gluing.

There is no volume mass at $S$, since its $\rho$-neighborhood has volume
$O(\rho^4)$. Volume is therefore regular Riemannian volume, and balls
have converging volumes. Indeed at almost every regular point away
from the center, the last regular segment of a minimizing path gives
$|\nabla d|=1$. Distance levels consequently have zero four-dimensional
measure. The regular limit has finite curvature energy by lower
semicontinuity on compact exhaustions. In $h_i^{-1}$ length units,
Equation~\eqref{geom:ricci-L3} gives a cubic Ricci integral
$O(h_i^2)$ on each fixed ball. Thus the normalized polar defect
vanishes, and Equations~\eqref{geom:almost-bishop} and
\eqref{geom:radial-laplacian} give exact Bishop monotonicity and
$\Delta(d^2/2)\leq4$ weakly on the regular part of the limit.

The finite exceptional set cannot separate positive-volume regular
components. Fix a regular initial point and a small regular endpoint
ball $U$ in a supposed other component, both a positive distance from
$S$. All approximating minimizing segments between them must enter
$N_\rho(S)$; otherwise a subsequential limit of such bounded segments
would connect the two components in the regular complement.
Their passage through $N_{C\rho}(S)$ has length at least $c\rho$,
at radii in a fixed compact subinterval of $(0,\infty)$.
The polar inequality gives, for a passage radius $s$ and endpoint
radius $r$,
\[
 P(r)\leq C P(s)+C\int_s^r m(u)P(u)\,du.
\]
Integrating over passage times, endpoint radii, and angles, and using
Fubini on these bounded ranges, gives
\begin{equation}\label{geom:nonseparation}
 c\rho\Vol(U)\leq
 C\left(\Vol(N_{C\rho}(S))+\int m\right).
\end{equation}
Cuts only decrease the polar integrals. First take the sequence
limit, eliminating $\int m$, and then let $\rho\downarrow0$.
The right side is $O(\rho^4)$, a contradiction. This argument also
applies to iterated tangent and end limits: a diagonal choice ensures
that the original physical radii still tend to zero.

Write $d$ for distance from a puncture, or from a fixed point when
studying infinity. At a puncture, finite energy makes the curvature
energy on annuli $\{r/4<d<4r\}$ tend to zero as $r\downarrow0$; at infinity the same
holds as $r\to\infty$. First, point-picking and the fixed energy
$e_0$ rule out unbounded $r^2|\Rm|$ on a smaller annulus. The inherited
local estimates then give $r^3|\nabla\Rm|\le C$ with further fixed
boundary margins. If $r^2|\Rm|\ge\varepsilon>0$ at a point of this
middle annulus, curvature remains at least $\varepsilon/(2r^2)$
on a ball of radius $c\varepsilon r$ contained in the wider annulus.
For each fixed $\varepsilon$, lower volume bounds give energy at
least $c\varepsilon^6$ on that ball, contradicting the vanishing
annular energy. Thus the scaled curvature tends to zero. Interior
stationary estimates, or interpolation with the inherited higher
bounds, give the same conclusion for every fixed number of scaled
derivatives. Tangent and end limits have flat regular annuli and only
the dilation center as a possible incomplete point. Their ball volumes
are exactly $vr^4$: the monotone ratio has a limit at the relevant end,
and dilation freezes that limit at every fixed radius.

Put $u=d^2/2$. The nonnegative distribution $4-\Delta u$ pairs to zero
with every compactly supported radial test by this volume formula
and $|\nabla d|=1$. Dominate an arbitrary nonnegative test by a radial
test on a larger compact annulus to see that the defect vanishes.
Elliptic regularity makes $u$ smooth and $\Delta u=4$. Bochner's formula
and $|\nabla u|^2=2u$ give $|\nabla^2u|^2=4$; since its trace is $4$,
equality in Cauchy's inequality gives $\nabla^2u=g$. Its gradient
curves produce the cone metric $dr^2+r^2g_L$. Flatness makes $L$ a
compact spherical manifold and nonseparation makes it connected.
Thus $L=S^3/\Gamma$ for a finite orthogonal group acting freely.
Two-sided volumes bound $|\Gamma|$. All fixed-ratio annuli sufficiently
far along the original end are therefore arbitrarily close, with
derivatives and radial distances, to round quotient annuli.
Enlarged overlapping annuli identify the groups.

\begin{lemma}[Curvature and chart improvement on necks]\label{geom:neck}
Suppose annular regions have inner and outer radii $r_-\ll r_+$ and
every escaping intermediate annulus converges smoothly to a round flat
quotient annulus. Assume for some $\sigma>0$ and all needed finite
orders that
\[
 r^2|\Ric|_{k+2,r}\leq
 C\big((r/r_+)^\sigma+(r_-/r)^\sigma\big).
\]
Fix $0<\mu<\min(1,\sigma)$. After removing fixed-factor boundary
collars, for this fixed exponent,
\begin{equation}\label{eq:neck-decay}
 r^2|\Rm|_{k,r}\leq
 C\big((r/r_+)^\mu+(r_-/r)^\mu\big).
\end{equation}
A single equivariant coordinate system on a shortened neck gives the
same bound for $g-\delta$, with scaled derivatives. Infinite
half-necks are allowed. One may also use slowly varying coefficients
on the boundary terms, when the forcing and boundary data obey their
corresponding bounds.
\end{lemma}

\begin{proof}
First choose the permitted fixed-factor boundary collars. For a fixed
finite derivative order, a fixed widened annular ratio, and a small
cone-chart tolerance, there are a fixed $K$ and a late index such that
all widened annuli with central radii in $[Kr_-,r_+/K]$ have these
smooth bounds. Otherwise choose increasingly late violating annuli
with radii in $[jr_-,r_+/j]$, $j\to\infty$. Both boundary ratios
escape, contradicting the assumed smooth convergence of the whole
annulus. Enlarge $K$ by a fixed factor to retain boundary margins.
The shortened boundaries $r'_-=Kr_-$ and $r'_+=r_+/K$ therefore have
uniformly bounded scaled curvature and the required derivatives.
For either exponent $\beta=\mu,\sigma$, their weights satisfy
\[
 (r/r'_+)^\beta+(r'_-/r)^\beta
   =K^\beta\big((r/r_+)^\beta+(r_-/r)^\beta\big).
\]
It is enough to prove the estimate on this shortened neck: these
fixed multipliers preserve both the forcing hypothesis and the
claimed conclusion. Relabel its boundary radii as $r_-,r_+$ below.

Write $w_\mu(r)=(r/r_+)^\mu+(r_-/r)^\mu$. If the bound fails, choose
approximate maximizing points of $r^2|\Rm|/w_\mu(r)$ with ratios
$Q_i\to\infty$, and call their radii $s_i$. Both boundaries escape
after dilation by $s_i^{-1}$: near either boundary the weight is bounded
below and scaled curvature is bounded. Put
$A_i=s_i^2|\Rm|(s_i)$. Intermediate annular flatness gives $A_i\to0$.
Divide the rescaled curvature by $A_i$. The maximum property and
\[
 \frac{w_\mu(rs_i)}{w_\mu(s_i)}\leq\max(r^\mu,r^{-\mu})
\]
give the divided tensor the bound
$Cr^{-2}\max(r^\mu,r^{-\mu})$.
The second Bianchi identity and commutation give
$\Delta\Rm=\Rm*\Rm+\nabla^2\Ric*$.
The divided quadratic term vanishes. The divided forcing vanishes on
fixed annuli since $w_\sigma(s_i)/A_i\leq C/Q_i$.
Interior elliptic estimates thus give a nonzero limiting Cartesian
tensor $H$ on the flat punctured cover, componentwise harmonic and
satisfying the uncontracted differential Bianchi identity.

Spherical harmonics have radial powers $j$ and $-j-2$,
$j=0,1,\ldots$. The two bounds with $\mu<1$ leave only
$H=r^{-2}C$ with a constant Cartesian tensor $C$. Bianchi becomes
\[
 x_iC_{jk ab}+x_jC_{ki ab}+x_kC_{ij ab}=0
       \qquad\text{for every }x\in\R^4.
\]
For fixed $a,b$ it says $x\wedge C_{\cdot\cdot ab}=0$ for every
covector $x$, forcing that two-form to vanish. This contradicts the
normalization. Derivatives follow by interior elliptic estimates.
For an infinite end use a remote artificial boundary and pass to its
limit. Slowly varying boundary coefficients obey the same ratio
estimate after decreasing $\mu$ to absorb their change on fixed
logarithmic intervals; the same divided forcing argument applies.

Here are the gauge and gluing details. On a fixed lifted round annulus,
with a larger annulus retained for margins, let
$S(V)=DV+(DV)^t$ and $K(p)$ be linearized curvature at $\delta$.
For $k\geq5$ there is a bounded linear extraction
$E:H^k(\operatorname{Sym}^2)\to H^{k+1}(T)$, modulo Killing fields,
with
\begin{equation}\label{geom:linear-gauge}
 E(SV)=V,\qquad
 \|(1-SE)p\|_{H^k}\leq C\|K(p)\|_{H^{k-2}}.
\end{equation}
Indeed regard the linearized connection
$C^a_{ij}=\tfrac12(\partial_i p_{ja}+\partial_jp_{ia}-\partial_ap_{ij})$
as a one-form in $i$, for each $a,j$. Its curl is $K(p)$.
The Neumann gradient projection writes it as $dU^a_j$ plus an
$H^{k-1}$ error controlled by $K(p)$ in $H^{k-2}$.
Symmetry in $i,j$ makes the curl of $U^a_jdx^j$ controlled; a second
projection writes $U^a=dV^a$ plus a controlled error.
Thus $C(p-SV)$ is controlled, and
$\partial_iq_{ja}=C(q)^a_{ij}+C(q)^j_{ia}$ controls $Dq$.
Absorb the remaining constant symmetric matrix by an affine field.

For clarity, the projection estimate is the div/curl estimate with
zero normal component. Integration by parts makes the boundary
second fundamental form an order-zero term and gives coercivity
modulo $L^2$. Compactness removes this term because a closed remainder
is a gradient on the simply connected annulus and the Neumann
projection makes it orthogonal to gradients. Higher estimates follow
by tangential differentiation in flattened boundary charts, recovering
normal derivatives from div and curl. Fix averages and Killing
ambiguity, or project onto the range of $S$, to obtain the stated
left inverse at all needed orders. Averaging over $\Gamma$ preserves
the estimates and gives equivariance.

For $g=\delta+p$ small in $H^k$, solve
$V=Ep-E((DV)^tDV)$ by contraction. Sobolev multiplication is bounded
since $k>4/2+1$. Put $q=p-SV-(DV)^tDV$, so $Eq=0$, and
$f=(\Id+V)^*\delta$. Since $f$ is flat, the difference of the nonlinear
curvature remainders of $g$ and $f$ has $H^{k-2}$ norm at most
$C\|p\|_{H^k}\|q\|_{H^k}$. Equation~\eqref{geom:linear-gauge} and
absorption give
\[
 \|q\|_{H^k}\leq C\|\Rm(g)\|_{H^{k-2}}.
\]
Changing coordinates by $\Id+V$ gives the local metric estimate.
Higher finite orders and Sobolev embedding give all required scaled
differentiability bounds.

Choose geometric radius steps and charts on much wider annuli whose
images contain whole distance bands with margins. On consecutive
overlaps, the transition satisfies
\[
 (D\Psi)^tD\Psi-I=O(\varepsilon_j),\qquad
 D^2\Psi=O(\varepsilon_j/r_j),\qquad
 \varepsilon_j=Cw_\mu(r_j).
\]
The second estimate follows from the Christoffel transformation law.
Integration gives
$\Psi=P_jx+b_j+O_{k,r_j}(\varepsilon_jr_j)$ for an affine isometry.
Comparing both charts in a larger converging cone chart, or using
nested bands in both directions, identifies their fundamental groups.
Average $P_j$ over this exact deck isomorphism and take its polar
factor to obtain an exact orthogonal intertwiner within the same
error. Project the translation to the invariant subspace; it is zero
when $\Gamma$ is nontrivial and free.

Align consecutive charts by these rigid motions. For the trivial
group, radial positioning bounds translation increments by $o(r_j)$.
Toward a puncture their sum converges; subtract its limit. At infinity
sum outward from an initial chart, giving $o(r_j)$ relative displacement
by geometric summation. Rotations do not affect radial position.
Interpolate the remaining errors on shorter overlaps; derivatives
of the interpolation cutoff cost $r_j^{-m}$ and preserve the scaled
$O(\varepsilon_j)$ bounds. Local invertibility follows from smallness.
For global injectivity, equal images have comparable radii and hence
lie in one common wide annular chart. There the glued map is $C^1$
close to the identity: close points are separated by derivative
bounds and points a fixed fraction of a radius apart by the $C^0$
bound. Radial margins ensure coverage of the shortened neck.
\end{proof}

\subsection{Removal of punctures and intrinsic singular points}
\label{geom:puncture-removal}

The neck estimate supplies coordinates in which a puncture has a
continuous metric completion. We now improve that completion to a smooth
metric on its local cover, and then prove that every actual concentration
point has a nontrivial quotient group. This second assertion is needed
both for organizing trees and for tracking them by their distances.
At a limit puncture Lemma~\ref{geom:neck}, with zero Ricci forcing,
gives on the cover
\begin{equation}\label{geom:puncture-decay}
 |D^k(g-\delta)|\leq C_kr^{\mu-k},\qquad \mu>0.
\end{equation}
Choose $0<\alpha<\min(\mu,1)$ and fill $g(0)=\delta$. Points separated
by a distance comparable to their radii are controlled by the
zeroth-order bound; closer points by the first-derivative bound along
their segment. Thus $g$ and $A^{ij}=\sqrt{\det g}\,g^{ij}$ are
$C^{0,\alpha}$, and after scaling $B_R$ to $B_1$,
$\|A(R\cdot)-I\|_{C^{0,\alpha}}=O(R^\mu)$.

Solve
\[
 \partial_i(A^{ij}\partial_jy^a)=0\quad\text{in }B_R,\qquad
 y^a=x^a\quad\text{on }\partial B_R.
\]
Coercivity gives the unique variational solution. For $w=y-x$,
\[
 \partial_i(A^{ij}\partial_jw^a)
       =-\partial_i(A^{ia}-\delta^{ia}),\qquad
 w^a|_{\partial B_R}=0.
\]
After rescaling to $B_1$, the forcing vector has
$C^{0,\alpha}$ norm $O(R^\mu)$. The weak-solution gradient estimate
\cite[Theorem~8.16 and Corollary~8.35]{GilbargTrudinger2001} gives
$\|Dy-I\|_\infty=O(R^\mu)$ and $y\in C^{1,\alpha}$. This estimate
also follows by absorbing the small coefficient perturbation in the
constant-coefficient divergence estimate. Thus $y$ is a coordinate
diffeomorphism for small $R$. Uniqueness makes it equivariant, and
$y(0)$ is $\Gamma$-fixed, so recentering is permitted. Put $B=Dy(0)$,
and then make a constant linear normalization. Since
$y(x)-y(0)-Bx=O(r^{1+\alpha})$, rescaled interior estimates after
subtracting the affine part give
$D^ky=O(r^{1+\alpha-k})$ for $k\geq2$. The forcing from that affine
part has the stronger coefficient exponent $\mu>\alpha$.
The inverse coordinates obey the same estimates. Consequently the
harmonic-coordinate metric satisfies
\begin{equation}\label{geom:harmonic-decay}
 G-G(0)=O(r^\alpha),\quad DG=O(r^{\alpha-1}),\quad
 D^2G=O(r^{\alpha-2}).
\end{equation}

Thus $G\in W^{2,p}$ across zero for all $p<4/(2-\alpha)$; take
$p_0=2+\alpha/2>2$. There is no distributional point defect:
integration by parts across $\partial B_\rho$ has first-derivative
boundary terms $O(\rho^3)$ and second-derivative terms involving
$DG$ of size $O(\rho^{2+\alpha})$. Both tend to zero.
The harmonic-coordinate Einstein equation holds almost everywhere:
\begin{equation}\label{geom:weak-einstein}
 G^{ab}\partial_a\partial_bG_{ij}=Q_{ij}(G^{-1},DG,DG).
\end{equation}
Here $Q$ is quadratic in $DG$ with smooth coefficients in $G^{-1}$;
this is the Lanczos harmonic-coordinate formula, as recorded in
\cite[Lemma~4.1, equation~(4.3)]{DeTurckKazdan1981}. For $2<p<4$, Sobolev gives
$DG\in L^{4p/(4-p)}$, so elliptic regularity improves the exponent by
\begin{equation}\label{geom:bootstrap-exponents}
 p_{k+1}=\frac{2p_k}{4-p_k}>p_k,\qquad
 \frac1{p_k}=\frac12-2^k\left(\frac12-\frac1{p_0}\right).
\end{equation}
After finitely many steps one can take $p>4$. If a step reaches
exactly $4$, the local embedding $W^{1,4}\subset L^q$ for every
finite $q$ permits the next exponent above $4$.

This gain does not assume the stronger regularity in advance.
On a sufficiently small ball the continuous leading coefficients
are uniformly close to a constant positive matrix. Localize $G$;
its equation has right side in the higher $L^q$ since the cutoff
terms contain only $G,DG$. The constant-coefficient $W^{2,q}$
inverse and a Neumann series for the small coefficient perturbation
give a $W^{2,q}$ solution. Uniqueness in the previously available
$W^{2,p}$ class identifies it with the localized metric.
Alternatively, the global Dirichlet regularity theorem
\cite[Theorem~4.2]{ChiarenzaFrascaLongo1993} gives this upgrade for the
localized metric components, which have zero boundary trace.
Keep the continuous leading coefficients equal to $G^{-1}$ near the
support of the cutoff and extend them to a constant positive matrix
outside a larger compact set. A convex interpolation
preserves ellipticity and gives bounded uniformly continuous, hence
vanishing-mean-oscillation, coefficients.
Once $p>4$, $G\in C^{1,\beta}$.
Equation~\eqref{geom:weak-einstein}, Schauder estimates, and
differentiation give smoothness. This proves the orbifold filling
in changed equivariant coordinates without an assumption about
Einstein deformation integrability.

A concentration point cannot have trivial group. If it did, its
smooth filling would have almost-Euclidean volume ratios on sufficiently
small balls. Smooth convergence on surrounding annuli, negligible
tip volume, and Equation~\eqref{geom:almost-bishop} propagate these
lower ratios to every smaller scale in the approximations: the total
error below original radius $h$ is $O(h^{1/2})$. Concentration and
point-picking give, at a much smaller scale, a complete nonflat
Ricci-flat limit with a regular nonflat region. Point-picked centers
are $o(h)$ from the smoothly filled center. Inclusions between their
radius-$h$ balls and the filled-center balls of radii $h\pm o(h)$
transfer the same lower starting ratios. Let the initial ratio approach
the Euclidean one, choosing the approximating index after each outer
radius; the summable comparison error tends to zero in physical units.
Every fixed ball of the smaller limit has at least
Euclidean volume ratio. At a regular center Bishop comparison gives
the opposite bound. Equality, by the radial Laplacian and Hessian
argument, forces a smooth flat cone from that center, a contradiction.
Thus each exceptional point is a genuine intrinsic singularity with
nontrivial group; in particular a metrically regular point of a limit
cannot conceal a further curvature concentration. The same argument
shows the end group of a nonflat complete limit is nontrivial: a trivial end gives Euclidean
asymptotic volume ratio, which is independent of the chosen center.

A complete flat limit here is a global Euclidean quotient with at
most one tip. Develop the universal cover of its regular locus.
The faithful free linear local holonomy at a puncture ensures that
its local fundamental group injects into that of the regular locus:
the composition of the inclusion with flat holonomy is the faithful
linear action of its quotient group. Hence its lifts contain full
punctured Euclidean balls. Fill their centers.
The resulting flat manifold is complete: a finite-length path
projects to a convergent path downstairs and eventually remains in
one lifted local ball. Adding isolated points preserves simple
connectivity in dimension four. Its development is therefore all
of Euclidean space. The deck group is finite. For any prescribed
finite number of distinct deck transformations, lifts of a regular
expanding ball and their transforms occupy that many sheets in a
comparably enlarged Euclidean ball. Every downstairs point has a
lift within comparable radius by path lifting; paths detour isolated
tips with arbitrarily small excess using their connected links.
Arbitrarily many sheets would force the downstairs Euclidean growth
ratio to zero, contrary to noncollapse. A finite Euclidean isometry
group fixes its barycenter and is orthogonal after translation.
Since singularities are isolated, its only possible singular point
is the origin. This finishes the small-scale orbifold, genuine-tip,
and end assertions of Theorem~\ref{thm:degeneration}. We can now use
the intrinsic singular points to identify the vertices and edges of
a finite tree.

\subsection{Finite bubble trees and the definition of the root scale}
\label{geom:trees-root}

We have shown that a flat limit has at most one genuine tip. This
excludes two separated concentrations without a larger nonflat ancestor.
Together with the slice energy bound it will produce a finite tree and
identify its outermost member by a continuous maximum.
The terminal space $Z$ has only finitely many concentration points.
Fix a sequence $t_i\uparrow T$ and write $\pi:M\to Z$ for the
terminal quotient. For each sufficiently large integer $m$, let
$E_{i,m}=\pi(\{l(\cdot,t_i)\leq2^{-m}\})$. Extract their Hausdorff
limits $F_m$ diagonally in the compact space $Z$, allowing an empty
limit if the sets eventually become empty. They are nested compact
sets. The bad-set cover and uniform convergence of distances show
that each $F_m$ is covered by at most $N$ balls of radius $C2^{-m}$,
where $N,C$ are independent of $m$. Consequently
$S=\bigcap_mF_m$ has at most $N$ points: $N+1$ distinct points
would violate such a cover for large $m$.

Let $K\subset Z\setminus S$ be compact. Nested compactness gives
an $m$ with $K\cap F_m=\varnothing$ and a fixed neighborhood of
$K$ disjoint from $F_m$. For all sufficiently large $i$, every
material point mapping to a slightly smaller neighborhood has
$l(\cdot,t_i)>r=2^{-m}$. Crucially, $r$ is now fixed. Choose one
such $i$ with $T-t_i<c r^2$. Forward spatial persistence on these
radius-$r$ balls controls the entire remaining interval $[t_i,T)$,
and a finite ball cover gives fixed material charts, bounded
derivatives, and smooth terminal convergence over $K$. Local
injectivity gives a single sheet. Thus the possible exceptional
set is fixed and finite throughout late time; this conclusion does
not infer uniformity from shrinking time windows. It completes the
terminal-convergence assertion of Theorem~\ref{thm:degeneration}.

If the exceptional set is empty, finitely many such charts bound
curvature globally. Their derivative estimates, metric evolution,
and compactness give a smooth limiting metric on the original
manifold at $T$. The DeTurck formulation, with
$W^a=g^{ij}(\Gamma^a_{ij}(g)-\Gamma^a_{ij}(g_{\rm fixed}))$,
has principal part $g^{ij}\partial_i\partial_jg$ after adding
$\mathcal L_Wg$. Short-time quasilinear parabolic existence and
pullback by $-W$ continue the flow on the same manifold. The smooth
limit and evolution equation give matching derivatives at $T$;
compare \cite[Theorem 14.1]{Hamilton1982} for the curvature continuation
criterion. Hence only the case of a genuine concentration point
needs further work. Fix a small neighborhood $\mathcal U$ in the
material manifold, pulled back from a $d_T$-ball, whose regular
boundary isolates one such point. Its punctured regular annuli have
finite limiting curvature energy.

For a sequence $t_i\uparrow T$, call two nonflat small-scale limits
in this cluster \emph{equivalent} if their scales have ratio bounded
above and below and their centers remain a bounded distance apart
in either scale. Pass to subsequences so all scale and separation
comparisons converge in the extended sense. Each nonflat limit
contains a regular ball with a uniformly positive energy. Indeed
its orbifold filling and end decay make its maximum curvature finite
and attained. Dilate so that maximum equals $1$. The complete
stationary curvature equation gives a universal derivative bound:
the differentiated-curvature product used above, now stationary,
satisfies $-\Delta Q\leq C-cQ^2$ for
$Q=(D+|\Rm|^2)|\nabla\Rm|^2$. Its maximum is controlled, using decay
at infinity and lifted maximum tests at quotient points.
Thus a definite ball around a curvature maximum has curvature
bounded below. Noncollapse, upper volume bounds, and the uniformly
finite number of tips allow removal of small tip neighborhoods
without exhausting that ball; some smaller regular ball remains
with definite curvature and volume. Its energy is at least a fixed
$e_1>0$.

Such regular energy balls belonging to inequivalent entries are
disjoint in the original slices for large $i$. Separated centers
give disjoint balls directly. If one scale is much smaller and
its center remains bounded on a larger scale, it cannot meet a
fixed regular ball there: smooth regular convergence bounds its
curvature radius from below proportionally to the larger scale.
The smaller entry must approach a tip. Therefore the total energy
bounds the number of inequivalent entries by $C/e_1$.

Construct a maximal list by adjoining, whenever possible on a
further subsequence, a nonflat limit inequivalent to those already
listed. Previous convergences and comparisons persist after
subsequence extraction. The energy bound forces this process to
terminate, including against tests on further subsequences.
The list is nonempty by high-curvature point-picking.
If two entries are separated by a distance much greater than both
scales, their separation-scale limit has two genuine tips. It cannot
be flat, since a flat limit has at most one tip. It is consequently a
new nonflat ancestor unless it is already in the list.
Comparable scales at bounded distance are equivalent; at separated
scales and bounded larger-scale distance the smaller entry approaches
a tip of the larger one. These observations define an ancestor order
on the finite list and show that it has a unique maximal entry.
Applying the same argument to entries at each tip gives a unique
outermost child in that branch, and every tip has a child by
concentration and point-picking. This is the finite rooted tree.

Every intermediate limit between consecutive entries is flat,
since a nonflat one would enlarge the maximal list. The child becomes
its cone tip, and the same conclusion holds between the root and
arbitrarily slowly vanishing outer radii. On these flat annuli
$l\geq cr$, so Equation~\eqref{eq:ricci-improvement} gives
$r^2|\Ric|_{k,r}\leq Cr$ in physical units. Lemma~\ref{geom:neck}
therefore supplies the finite necks and positive decay exponents.
This completes the tree assertion of Theorem~\ref{thm:degeneration}.

These convergences identify actual slice regions with the core/neck
domains of Section~\ref{sec:static-input}. Use compact regular embeddings
on truncated vertices and the single annular coordinates on each
joining region. Dividing physical coordinates by the corresponding
vertex scale gives its end coordinates. If the parent and child
scales are $A$ and $Ae^{-2L}$, their end variables satisfy
\[
 \tau_p=\log(A/r),\qquad
 \tau_c=\log(r/(Ae^{-2L})),\qquad \tau_p+\tau_c=2L.
\]
On the parent half $\tau_p\leq L$,
$e^{-\mu\tau_p}+e^{-\mu\tau_c}\leq2e^{-\mu\tau_p}$;
the child half has the analogous bound. Thus the two-sided neck
estimate supplies exactly the one-sided half-end weights used there.
Fixed initial end collars are absorbed by enlarging constants.

On fixed overlaps the annular charts converge, after a subsequence,
with all required derivatives; they prescribe end charts for the
limiting vertices by diagonal extraction. Adjust the regular
embeddings by the small resulting overlap differences. Regular
convergence covers each vertex except holes about its tips;
annular convergence covers the intervening distances to its children,
and at a leaf a bounded rescaled core is entirely regular.
The same construction passes through a fixed compact root collar if
the exterior is later replaced. It is a sequential construction:
there is no assertion that one tree or one family of identifications
is smooth in time. It gives only a positive joining exponent.
The stronger exponent greater than one on the unjoined exterior
will be supplied by Theorem~\ref{thm:collar-extension} before
Theorem~\ref{thm:tree-inequality} is applied.

\begin{proposition}[A continuous outermost scale]\label{prop:root-scale}
There is a fixed small $c_*>0$ such that the following definition
detects the unique outermost entry. Choose a bounded continuous
$\psi:[0,1]\to[0,\infty)$, zero on $[0,c_*]$ and strictly positive
on $[2c_*,1]$, and set
\begin{equation}\label{geom:root-definition}
 a(t)=\sup_{y\in\mathcal U}
 l(y,t)\psi\big(l(y,t)^2|\Rm|(y,t)\big).
\end{equation}
After shrinking $\mathcal U$, this supremum is attained in its
interior for every late time; it is positive and continuous, and
$a(t)\to0$. For every sequence of late times, scale $a(t)$ and any
maximizing point represent the outermost nonflat entry. Every other
nonflat entry has scale bounded by a constant times $a(t)$.
The constants may be fixed after extracting the sequential tree.
\end{proposition}

\begin{proof}
The normalized regular energy balls just constructed provide uniform
detection. In the normalization $\max|\Rm|=1$, the derivative bound
keeps $|\Rm|\ge1/2$ on a ball of fixed radius. Removing small
neighborhoods of the uniformly bounded number of tips leaves a point
a fixed positive distance from every tip, by the volume bounds.
Its curvature radius is bounded below: a sufficiently small fixed
ball is regular and $|\Rm|\le1$ there. It is bounded above by the
nonzero curvature at its center. Hence at this point
$l^2|\Rm|\geq2c_*$ for a fixed small $c_*$, and $l$ is comparable
to the vertex scale. This assertion transfers to approximating slices.
Here curvature radius in a limit is stopped by its exceptional
points: a ball containing such a point has unbounded approximating
curvature. On compact regular balls, smooth convergence verifies
the opposite inequality in the radius definition.

On the punctured terminal neighborhood the quantity $l^2|\Rm|$
tends uniformly below $c_*/2$ toward the puncture. Otherwise regular
points approaching it with $l^2|\Rm|\geq c_*/2$ yield disjoint
definite-energy balls, using curvature derivatives and noncollapse.
Their radii are at most a constant times distance to the puncture,
since concentration stops the curvature radius. Taking the
terminal limit first keeps these balls regular; their definite
energies contradict finite energy arbitrarily near the puncture.
Shrink $\mathcal U$ within this region. On every compact subset of
its punctured terminal regular part, $l^2|\Rm|<c_*$ eventually.
In particular the integrand in Equation~\eqref{geom:root-definition}
vanishes near its boundary. It is continuous on a fixed compact
interior set, so its maximum is attained and depends continuously
on time. Detection proves positivity. Every sequence of maximizing
points converges to the concentration point and has $l\to0$;
hence $a(t)\to0$.

At a maximizer, $l\geq a/\|\psi\|_\infty$ and
$l^2|\Rm|>c_*$. Its $l$-scale limit is therefore nonflat. If
$l/a\to\infty$, detection at another regular point of that same limit
would give a value of $l\psi$ bounded below by a positive multiple
of $l$, contradicting the definition of $a$. Thus $l\asymp a$ at
maximizers. Conversely detection on any listed nonflat limit bounds
its scale by $Ca$. The maximal list contains the maximizing
$a$-scale limit, and its unique outer ancestor cannot have larger
order than $a$. It is consequently the same entry. This is precisely
the outermost scale; a deepest bubble or the largest curvature may
occur at a much smaller scale.
\end{proof}

\subsection{Motion on short intervals}

The scale just defined must be compared at different times on the same
material manifold. We first use the genuine quotient tips to control
curvature radii along minimizing segments, then construct dense sets of
material points whose distance changes are small.
\begin{lemma}[Curvature radius along minimizing segments]\label{geom:geodesic-radius}
For any minimizing segment $\gamma:[0,D]\to M$ at a late time,
parametrized by arclength,
\[
 l(\gamma(s),t)\geq c\min(l_0,s,D-s),\qquad 0<s<D.
\]
Consequently, whenever $d_t(x,y)\leq Ch$ with $h$ small, its upper
and lower time derivatives satisfy
\begin{equation}\label{geom:distance-log}
 |\partial_t d_t(x,y)|\leq
 C'\left(1+\log^+(h/l(x,t))+\log^+(h/l(y,t))\right).
\end{equation}
Time derivatives can be interpreted almost everywhere or by barriers.
\end{lemma}

\begin{proof}
If the bound fails, choose minimizing segments and interior points
$z_i=\gamma_i(s_i)$ such that, with
\[
 r_i=l(z_i,t_i),\qquad L_i=\min(l_0,s_i,D_i-s_i),
 \qquad r_i/L_i\longrightarrow0,
\]
the scale $h_i=\sqrt{r_iL_i}$ satisfies
$h_i\to0$, $r_i/h_i\to0$, and $h_i\le L_i$.
Restrict to the centered subsegments of length $2h_i$ and rescale
by $h_i^{-1}$. Their limits are minimizing segments with endpoints
at distance one on both sides of the basepoint. Since $l(z_i,t_i)/h_i\to0$,
that basepoint is a genuine exceptional point of the slice limit.
Its tangent cone has link $S^3/\Gamma$ with $\Gamma$ nontrivial and free.
This link has diameter strictly less than $\pi$: if two orbits were
distance $\pi$ apart, every translate of one representative would
have to be the antipode of the other, forcing its orbit to be a
singleton, contrary to freeness. Blowing up the limiting segment at
this point therefore gives a minimizing segment through the cone
vertex with positive length on each side. Joining nearby points by
the shorter cone chord contradicts minimality.

For the derivative bound use first variation
$|\partial_t d|\leq\int_\gamma|\Ric|\,ds$ and
Equation~\eqref{eq:ricci-improvement}. Near $x$ the spatial
Lipschitz property also gives $l(\gamma(s))\geq l(x)-s$.
Combining this with the interior bound gives
$l(\gamma(s))\geq c(l(x)+s)$ on the first half of the segment,
up to the harmless cap, and the analogous expression on the other
half. Integration of their reciprocals gives
Equation~\eqref{geom:distance-log}.
\end{proof}

\begin{lemma}[Uniform short motion]\label{geom:slow-motion}
Let $h_i\to0$ and let $J_i$ be time intervals of length
$\delta_i=o(h_i)$ tending to $T$. Distances initially of order
$O(h_i)$ change by $o(h_i)$ throughout $J_i$. If a material point
has $l\geq c h_i$ at one time, it has $l\geq c' h_i$ throughout
(after extraction for a sequential assertion). Regular compact
charts at this scale persist, and their scaled metric coefficients,
with any fixed finite number of scaled derivatives, change by
$O(\delta_i/h_i)$ in fixed material coordinates.
\end{lemma}

\begin{proof}
Put $q_i=\delta_i/h_i\to0$ and, on material points, define
\[
 I_i(x)=\int_{J_i}\log^+(h_i/l(x,t))\,dt.
\]
The distribution formula and Equation~\eqref{geom:bad-volume} give
\[
 \int_M\log^+(h_i/l)\,d\mu_t
 =\int_0^\infty\Vol_t\{l<h_ie^{-v}\}\,dv\leq Ch_i^4.
\]
Volume forms are comparable across $J_i$, so Fubini gives
$\int I_i\leq C\delta_i h_i^4$. Discard material points with
$I_i>h_i\sqrt{q_i}$. Their volume at every time is at most
$Ch_i^4\sqrt{q_i}$. Lower ball volumes therefore make the retained
set $O(h_iq_i^{1/8})$-dense simultaneously at all times.
For two retained endpoints in a bounded rescaled distance range,
Equation~\eqref{geom:distance-log} integrates to
\[
 |d_t(x,y)-d_s(x,y)|
 \leq C(\delta_i+h_i\sqrt{q_i})=o(h_i).
\]
A first-crossing argument keeps initially bounded distance ranges
bounded. Its reverse version prevents initially remote retained
points from entering a fixed bounded range.

Choose finite increasingly dense labeled nets of retained material
points on bounded balls. Their pairwise distance matrices remain
unchanged in the limit, uniformly in time. Every temporal
subsequence thus has the same pointed metric limit under these
labels. Exceptional points in all such limits are intrinsic genuine
singularities, and each metrically regular compact set has smooth
regular convergence, by the exclusion of trivial-group concentration.

Now take a material point regular at one time, but whose minimum
$l/h_i$ over $J_i$ supposedly tends to zero. By reversing time
direction in the distance comparison if needed, stop at its first
crossing of $l/h_i=\eta_i$, where $\eta_i\downarrow0$ is chosen
above that attained minimum and sufficiently slowly that
$q_i\log(1/\eta_i)\to0$. Until the stopping time its integrated
endpoint logarithm is at most
$\delta_i\log(1/\eta_i)=o(h_i)$. Its distances to retained anchors
therefore stay unchanged to $o(h_i)$. At the crossing it is a
singular point in the common pointed limit with the same anchor
distance vector as its initially regular point, a contradiction.
This proves the asserted uniform positive lower radius.

After taking a subsequence, the other alternative for a tracked point
is $\sup_{J_i}l/h_i\to0$. Such a point initially near a good anchor
must stay near the same singular point. Choose a large anchor radius
different from the finitely many singular-point radii in the common
limit. Escape to that sphere would give a singular point there by
subsequential extraction, which is impossible. In the resulting
bounded ball choose finitely many good labels whose distance vectors
distinguish the singular points. At arbitrary intermediate times,
the persistently bad point's vector is uniformly close to one of
these finitely many distinct vectors; otherwise take a contradicting
temporal subsequence. Its vector is continuous in time, so it stays
close to the same vector throughout. Increasing the finite anchor
nets identifies all its limiting positions with the same tip.

These alternatives prove tracking for every point: any sequence
violating it has either a uniformly regular time subsequence, handled
by the first argument, or a persistently bad subsequence, handled by
the second. Two tracked points can be placed in a common bounded
good-anchor range, so their distance changes by $o(h_i)$.
Finally, on persisting regular charts,
$|\nabla^k\Ric|\leq C_k h_i^{-1-k}$. Integrating the metric and
connection evolution in the fixed initial chart gives
$O(\delta_i/h_i)$ change of the scaled metric and its derivatives.
The coordinate conversion is an induction on derivative order,
using the already controlled lower-order connection terms.
\end{proof}

\begin{lemma}[Time required for a fixed scale change]
\label{geom:scale-change-time}
Fix $u>0$. Let $s_i<t_i<T$ tend to $T$, where $t_i$ is the first time
after $s_i$ for which $|\log(a(t_i)/a(s_i))|=u$. Then
\[
 \frac{t_i-s_i}{a(s_i)^2}\longrightarrow\infty.
\]
\end{lemma}

\begin{proof}
Otherwise a subsequence of these intervals has duration $O(a^2)$,
where $a=a(s_i)$. Rescale lengths by $a$. Their duration is $o(a)$
in physical units, so
Lemma~\ref{geom:slow-motion} freezes the pointed root metric and its
regular curvature in the limit. The radius in the definition of
$a$ converges on regular parts to the radius computed in that
limit with its tips as stops. For the lower bound use smooth
convergence on a slightly smaller admissible ball, whose curvature
threshold has a strict margin. For the upper
bound a ball exceeding the limiting radius either contains a
curvature violation on a regular compact set or encounters a
genuine concentration tip; in either case it is inadmissible in
the approximations.

Maximizers at both times are regular on this scale by
Proposition~\ref{prop:root-scale}. They remain in a common bounded
root range: a nonflat patch at the first time persists at the other,
and uniqueness makes all comparable nonflat scales the same tree
entry. Thus the rescaled supremum in
Equation~\eqref{geom:root-definition} is computed on regular compact
regions in one common pointed limit and has the same value at both
times. A fixed-factor change is impossible.
\end{proof}

\begin{proposition}[Buffered scale-change intervals]\label{prop:buffered-intervals}
If a concentration point exists, there are parabolically rescaled
intervals $I=[-\theta,1+\theta]$, for one fixed $\theta>0$, with
physical squared length units $d_i\to0$, such that the root scale
throughout $I$ is comparable to a number $a_i\to0$ and changes by a
fixed factor between $0$ and $1$. Put $\epsilon_i=\sqrt{d_i}$.
In these units
\[
 |R|\leq C\epsilon_i^2,\qquad
 |\Ric|_{k,r}\leq C_k\epsilon_i/r
\]
at regularity scales $r$ in a bounded range. For $a_i\ll r\lesssim1$
the cluster has flat quotient annular limits with nontrivial group.
At any fixed time, allowable root centers differ by $O(a_i)$.
\end{proposition}

\begin{proof}
Fix $u>0$ and choose successive late times $t_j$ at which
$\log a$ first changes by $u$ in absolute value from its value at
$t_{j-1}$. Positivity, continuity, and $a(t)\to0$ ensure infinitely
many such times with $t_j\uparrow T$. Put
$d_j=t_{j+1}-t_j$ and $a_j=a(t_j)$.
Lemma~\ref{geom:scale-change-time} gives $d_j/a_j^2\to\infty$.
Also $a_j\geq a_0e^{-uj}$, so eventually
\[
 d_j\geq c e^{-2uj}.
\]
Choose $0<\theta<\min(1,e^{-2u})$. Arbitrarily late indices satisfy
$d_{j-1},d_{j+1}\geq\theta d_j$. Otherwise every sufficiently late
index has a neighbor smaller by a factor $\theta$. Starting so late
that its duration is smaller than all durations in the finite
excluded initial segment, follow such neighbors. This chain cannot
enter that initial segment and cannot reverse direction: reversal
would return to the immediately preceding, strictly larger term.
It must run rightward forever. It then gives
$d_{j+k}\leq\theta^k d_j$, contradicting the displayed exponential
lower bound because $\theta<e^{-2u}$.

Use these central steps as the unit time interval. Their neighboring
steps provide buffers of size $\theta$. The first-change definition
makes $a(t)$ comparable to $a_j$ on the central step and both
buffers, with constants at most fixed powers of $e^u$.
In length units $\sqrt{d_j}$ its value is
$a_i=a_j/\sqrt{d_j}\to0$. The central endpoint factor is $e^{u}$
or $e^{-u}$; pass to a subsequence if its sign matters.
The curvature estimates rescale to the claimed bounds since
$|R_{\rm new}|=d_j|R|$ and
$|\Ric|_{k,r,\rm new}\leq C_k\sqrt{d_j}/r$.
Intermediate annuli are flat by maximality of the bubble list.
Their tip is genuine, so their quotient group is nontrivial.
The center bound follows because all maximizing points represent
the unique outermost entry.
\end{proof}

\begin{figure}[htbp]
 \centering
 \begin{tikzpicture}[
  x=1cm,y=1cm,>=Stealth,
  every node/.style={font=\small},
  central/.style={draw=blue!55!black,line width=1.2pt},
  buffer/.style={draw=orange!75!black,line width=1.2pt}]
  \node[anchor=west] at (0,1.0) {physical time};
  \draw[buffer] (0,0)--(2,0);
  \draw[central] (2,0)--(9,0);
  \draw[buffer] (9,0)--(11,0);
  \draw[->] (11,0)--(11.7,0) node[right] {$t$};
  \foreach \x in {0,2,9,11} \draw (\x,-.1)--(\x,.1);
  \node[above] at (0,.13) {$t_{j_i}-\theta d_i$};
  \node[above] at (2,.13) {$t_{j_i}$};
  \node[above] at (9,.13) {$t_{j_i+1}$};
  \node[above] at (11,.13) {$t_{j_i+1}+\theta d_i$};
  \node[below,text=orange!65!black] at (1,-.13) {$\theta d_i$};
  \node[below,text=blue!55!black] at (5.5,-.13) {$d_i$};
  \node[below,text=orange!65!black] at (10,-.13) {$\theta d_i$};
  \draw[->] (5.5,-.65)--(5.5,-1.55)
    node[midway,right=6pt] {$s=(t-t_{j_i})/d_i$};
  \node[anchor=west] at (0,-1.55) {interval units};
  \draw[buffer] (0,-2)--(2,-2);
  \draw[central] (2,-2)--(9,-2);
  \draw[buffer] (9,-2)--(11,-2);
  \draw[->] (11,-2)--(11.7,-2) node[right] {$s$};
  \foreach \x in {0,2,9,11} \draw (\x,-2.1)--(\x,-1.9);
  \node[below] at (0,-2.13) {$-\theta$};
  \node[below] at (2,-2.13) {$0$};
  \node[below] at (9,-2.13) {$1$};
  \node[below] at (11,-2.13) {$1+\theta$};
  \node at (5.5,-3.05)
    {$g_i(s)=d_i^{-1}g(t_{j_i}+d_i s),\qquad
      \epsilon_i=\sqrt{d_i},\qquad
      a_i=\dfrac{a(t_{j_i})}{\epsilon_i}\longrightarrow0$};
\end{tikzpicture}

 \caption{Parabolic normalization of a selected buffered interval.
Write $j=j_i$ for the selected indices and
$d_i=t_{j_i+1}-t_{j_i}$.  The neighboring intervals each have duration
at least $\theta d_i$; only the shortened buffers are drawn.
The root scale is comparable to $a_i$ throughout the normalized
interval, and its values at $0$ and $1$ differ by the fixed factor
$e^u$ or $e^{-u}$.  The diagram is schematic and does not assert
monotonicity of the root scale.}
 \label{fig:buffered-intervals}
\end{figure}

\begin{corollary}[Motion in interval units]\label{geom:rescaled-motion}
On the intervals of Proposition~\ref{prop:buffered-intervals},
geometry at length $h\lesssim1$ moves negligibly over windows
$\Delta t$ satisfying $\epsilon_i\Delta t/h\to0$.
Regular chart coefficients change by
$O(\epsilon_i\Delta t/h)$ with every fixed number of scaled
derivatives. At $h\asymp a_i$, old and new centers are $O(a_i)$
apart; at $a_i\ll h\lesssim1$ they are $o(h)$ apart.
\end{corollary}

\begin{proof}
The corresponding physical duration divided by physical length is
$d_i\Delta t/(\sqrt{d_i}h)=\epsilon_i\Delta t/h$, so
Lemma~\ref{geom:slow-motion} applies. At root scale a nonflat regular
patch persists, and uniqueness identifies the corresponding entry
at both times, giving bounded center displacement. At larger
intermediate scale the root is the intrinsic cone tip and its
location persists by the bad-point tracking argument. Another
concentration at distance comparable to $h$ in the same shrinking
terminal cluster would give two tips at their separation scale,
hence a nonflat ancestor larger than the outermost entry, a
contradiction. This proves the $o(h)$ displacement.
\end{proof}

\begin{proposition}[Uniform exterior geometry]
\label{geom:exterior-geometry}
Fix a sufficiently high finite derivative order $k_{\mathrm{geo}}$
and numbers $0<\mu<\sigma<1$. There are fixed constants
$N_{\mathrm{geom}},C,c_{\mathrm{geom}}>0$ and outer radii
$R_i^{\mathrm{out}}\to\infty$ in interval units, with
$\epsilon_iR_i^{\mathrm{out}}\to0$, with the following property.
At every time in the buffered interval the exterior region admits
one equivariant coordinate system whose annuli satisfy
\begin{equation}\label{geom:uniform-exterior}
 r^2|\Rm|_{k_{\mathrm{geo}},r}+|g-\delta|_{k_{\mathrm{geo}},r}
 \le C\Theta_i(r),\qquad
 \Theta_i(r)=(a_i/r)^\mu+\delta_i r^\mu,\qquad
 \delta_i=(R_i^{\mathrm{out}})^{-\mu}\longrightarrow0
\end{equation}
for $N_{\mathrm{geom}}a_i\le r\le c_{\mathrm{geom}}R_i^{\mathrm{out}}$.
The charts cover whole shortened distance bands with boundary margins;
their nontrivial quotient group is fixed after subsequence extraction.
Moreover $\epsilon_i=o(\delta_i)$. The constant $C$ is uniform in
time and radius and is unchanged when the lower bound is restricted
to $Na_i$ for any larger fixed $N$.
\end{proposition}

\begin{proof}
Choose actual outer radii $R_i^{\mathrm{out}}\to\infty$ in interval
units with $\epsilon_iR_i^{\mathrm{out}}\to0$, slowing their growth
if necessary. For arbitrary choices $t_i\in I$ and
$r_i/a_i\to\infty$ with $\epsilon_i r_i\to0$, the maximal-tree
and root-detection arguments force the annular limit to be a flat
quotient cone. Otherwise a nonflat limit at scale $r_i$ would, by
detection on a regular patch, give $a_i\geq c r_i$. This sequential
assertion implies that, for any fixed small cone-chart tolerance,
one fixed $N_{\mathrm{geom}}$ makes all widened annuli
$N_{\mathrm{geom}}a_i\leq r\leq R_i^{\mathrm{out}}$ sufficiently close to flat,
uniformly in $t\in I$ for large $i$. Failure would choose
$N_i\to\infty$ and violating times and radii, contradicting the
same assertion.

On these annuli $l\geq cr$, and hence
\[
 r^2|\Ric|_{k_{\mathrm{geo}}+2,r}\leq C\epsilon_i r
 \leq C\epsilon_iR_i^{\mathrm{out}}
             (r/R_i^{\mathrm{out}})^\sigma.
\]
Apply the ordinary two-boundary proof of Lemma~\ref{geom:neck}
with $r_-=N_{\mathrm{geom}}a_i$, $r_+=R_i^{\mathrm{out}}$, and the fixed exponent
$\mu$. Its constant is uniform in time: a violation chooses
maximizing pairs $(t_i,r_i)$; both boundary radii escape under
normalization, the annular metric tends to the flat cone, and the
divided Ricci forcing is at most $C/Q_i$. The harmonic-tensor and
Bianchi contradiction uses only $\mu<1$ and the flat cover,
independently of an extracted nonflat tree. This proves
Equation~\eqref{geom:uniform-exterior} on shortened
annuli, with all required derivatives. Fixed boundary factors are
included in $N_{\mathrm{geom}}$ and $c_{\mathrm{geom}}$.
Since $\epsilon_iR_i^{\mathrm{out}}\to0$ and $\mu<1$, one has
$\delta_i/\epsilon_i\to\infty$. Fixed wider collars merely
enlarge $N_{\mathrm{geom}}$ and $C$ once. These geometric choices
are fixed before the later analytical collar parameter $N$ is chosen.
Increasing $N\geq N_{\mathrm{geom}}$ only restricts this existing
estimate to $r\geq Na_i$ and preserves its constant; it does not
reconstruct the neck estimate with a new constant depending on $N$.
At scale one the regular charts persist
through the entire buffered interval by
Corollary~\ref{geom:rescaled-motion}; their quotient group is
therefore fixed after subsequence extraction. Widened overlapping
annuli identify that group and preserve the distance-band margins
in the chart construction at each time. This proves the exterior
estimate uniformly in all times and radii. Only the joining-end
exponents of individual extracted trees elsewhere in the argument
may be decreased after extraction.
\end{proof}
